\section{Background}
We shall work over a field $k$ in this article. We shall often need to deal with tensor products of $k$-vector spaces, and we shall
abbreviate $\otimes_k$ to $\otimes$.

If $G$ is a group and $k$ is a field, then we shall write $kG$ for the group algebra of $G$ over $k$. By a \textit{$kG$-module}, we shall
mean a right $kG$-module of finite $k$-dimension.
If $G$ is a group with a subgroup $H$, then for a field $k$ we shall write the operations of induction and restriction of modules
between the group algebras $kG$ and $kH$ as $\uparrow^{G}_{H}$ and $\downarrow^{G}_{H}$, with the field being implicit.

Mackey's theorem is a fundamental result in finite group theory which describes the interaction of the operations of induction and
restriction. If $G$ is a finite group and $H$ is a subgroup of $G$, then for $g \in G$ we define $H^g$ to be the subgroup
$\{g^{-1}hg \mid h \in H\}$ of $G$, and we call this the \textit{conjugate subgroup of $H$ by $g$} (note that $H^g$ is isomorphic to $H$).
Further, if $X$ is a $kH$-module, then we define $X^g$ to be the $kH^g$-module with underlying vector space $X$ and action given by
$x(g^{-1}hg) = xh$ for $x \in X$ and $h \in H$. We call this the \textit{conjugate module of $X$ by $g$}.

\begin{theo}\label{mackey:thm}\textup{(}Mackey's Theorem \cite[Theorem 3.3.4]{BRCI}\textup{)}
Let $G$ be a finite group with subgroups $H$ and $K$, let $\mathcal{U}$ be a complete non-redundant system of $(H,K)$-double coset
representatives in $G$, and let $X$ be a right $kH$-module. Then we have a decomposition of right $kK$-modules
\[
X\!\uparrow^{G}_{H}\downarrow^{G}_{K} \;\;\cong\;
\bigoplus_{u \in \mathcal{U}} X^u\downarrow^{H^u}_{H^u\cap K}\uparrow^{K}_{H^u\cap K}.
\]
\end{theo}

\subsection{Filtrations and series}\label{filt_ser:subsec}
Let $kG$ be a group algebra over a field, let $M$ be a $kG$-module and $X_1,\ldots,X_t$ also be $kG$-modules. A \textit{series for $M$
over $X_1,\ldots,X_t$} is a series of submodules
\begin{equation}\label{M_series:eq}
  M = M_n \supseteq M_{n-1} \supseteq M_{n-2} \supseteq \cdots \supseteq M_1 \supseteq M_0 = 0
\end{equation}
such that each quotient $\frac{M_l}{M_{l-1}}$ is isomorphic to some $X_i$. If such a series exists, we say that $M$ has a
\textit{filtration by the modules $X_1,\ldots,X_t$}. If $\frac{M_1}{M_0} = M_1$ is isomorphic to $X_l$, then we say that $X_l$ occurs at
the \textit{bottom} of the filtration or series. Further, let $f:\{1,\ldots,n\} \longrightarrow \{1,\ldots,t\}$ be any function such
that for each $l$, $X_{f(l)}$ is isomorphic to the quotient of $M_l$ by $M_{l-1}$, and let $\alpha_i = |f^{-1}(i)|$. We then say that
\eqref{M_series:eq}
is a \textit{series for $M$ over $X_1,\ldots,X_t$ admitting the multiplicities $(\alpha_i)_i$}, and that $M$ has a
\textit{filtration by the modules $X_1,\ldots,X_t$ admitting the multiplicities $(\alpha_i)_i$}.
Note that we have not assumed that the modules $X_i$ are pairwise non-isomorphic. If there are isomorphisms between the modules
$X_i$, then the multiplicities in a filtration are \textit{not uniquely determined} by the series of submodules, and so the same series
of submodules will admit different multiplicities. Even if the modules $X_i$ are pairwise non-isomorphic, so that the multiplicities
\textit{are} uniquely determined by the series, the multiplicities are \textit{not} in general uniquely determined by the module,
as the same module can have two series of submodules where the $X_i$ occur with different multiplicities.

\subsection{Combinatorics}
We now review a few combinatorial concepts. We assume that the reader is already familiar with these notions, and so our treatment will
be brief.

Recall that a \textit{composition} of $n$ is a tuple of non-negative integers adding up to $n$. We call $n$ the \textit{size} of $\alpha$
